% Lesefassung zum ausgelagerten Beweisgang über induzierte Potenzisomorphismen.
\section*{Von Elementen zu Teilmengen}
\hypertarget{power-isomorphisms.reading}{}

Ein Isomorphismus benennt die Elemente einer Halbgruppe um und erhält
 dabei ihre Multiplikation. Für eine nichtleere Teilmenge liegt dieselbe
 Vorschrift nahe: Wir benennen jedes ihrer Elemente um. Warum entsteht
 so wieder ein Isomorphismus? Wir zeigen zuerst die Erhaltung des
 Mengenprodukts und anschließend die Bijektivität der Bildabbildung.
 Die formalen Ableitungen stehen in den \PowerIsoProofsLink{};
 der abschließende Satz bleibt in Band~28 (\PowerIsoMainLink{}).

Seien \((A,\star)\) und \((B,\diamond)\) Halbgruppen. Ihre
nichtleeren Teilmengen bilden die Mengen
\[
  \PowNE{A}=\{X\subseteq A\mid X\ne\varnothing\},
  \qquad
  \PowNE{B}=\{U\subseteq B\mid U\ne\varnothing\}.
\]
Das \emph{Mengenprodukt} enthält sämtliche Produkte von je einem
Element aus dem linken und dem rechten Faktor:
\[
  \begin{aligned}
    X\star_{\mathcal P}Y
      &=\{x\star y\mid x\in X,\ y\in Y\},\\
    U\diamond_{\mathcal P}V
      &=\{u\diamond v\mid u\in U,\ v\in V\}.
  \end{aligned}
\]
Dies ist \FormulaRefAuto{PowerSemigroupProductDef}[def]. Mehrfach
entstehende Werte werden nur einmal aufgenommen. Die Reihenfolge
der Faktoren bleibt fest; Kommutativität wird nicht vorausgesetzt.
Das Mengenprodukt ist wieder nichtleer, und die Assoziativität
überträgt sich von der Grundoperation. Nach
\FormulaRefAuto{NonemptyPowerSemigroup}[thm] sind deshalb auch
\((\PowNE{A},\star_{\mathcal P})\) und
\((\PowNE{B},\diamond_{\mathcal P})\) Halbgruppen.

Ein \emph{Homomorphismus} \(f\colon A\to B\) erfüllt
\[
  f(x\star y)=f(x)\diamond f(y)\qquad(x,y\in A).
\]
Ist er bijektiv, heißt er \emph{Isomorphismus}
(\FormulaRefAuto{SemigroupHomDef}[def] und
\FormulaRefAuto{SemigroupIsoDef}[def]). Für \(X\subseteq A\) schreiben wir
\[
  f[X]=\{f(x)\mid x\in X\}.
\]
Während \(f(x)\) den Wert an einem Element bezeichnet, ist
\(f[X]\) die Menge aller Werte an Elementen von \(X\).
Diese Bildmenge wird zum Wert der induzierten Potenzabbildung.

\clearpage
\section*{Ein Homomorphismus erhält das Mengenprodukt}
\hypertarget{power-isomorphisms.reading.product}{}

Sei nun \(f\colon A\to B\) ein Halbgruppenhomomorphismus, und seien
\(X,Y\in\PowNE{A}\). Die entscheidende Gleichheit lautet
\[
  \boxed{f[X\star_{\mathcal P}Y]
    =f[X]\diamond_{\mathcal P}f[Y].}
\]
Links bilden wir zunächst alle Produkte in \(A\) und wenden danach
\(f\) auf die Ergebnisse an. Rechts bilden wir zunächst die beiden
Bildmengen in \(B\) und multiplizieren anschließend deren Elemente.
Um die Gleichheit zu beweisen, prüfen wir die beiden Inklusionen.

\medskip\noindent\textbf{Vom Bild des Produkts zum Produkt der Bilder.}
Sei \(z\in f[X\star_{\mathcal P}Y]\). Nach der Bedeutung einer
Bildmenge gibt es ein \(w\in X\star_{\mathcal P}Y\) mit \(z=f(w)\).
Dass \(w\) zum Mengenprodukt gehört, bedeutet wiederum: Es gibt
\(x\in X\) und \(y\in Y\), für die \(w=x\star y\) gilt.
Die Homomorphieeigenschaft liefert
\[
  z=f(w)=f(x\star y)=f(x)\diamond f(y).
\]
Wegen \(x\in X\) und \(y\in Y\) gehören die beiden Faktoren zu
\(f[X]\) beziehungsweise \(f[Y]\). Ihr Produkt gehört also zu
\(f[X]\diamond_{\mathcal P}f[Y]\). Damit ist gezeigt:
\[
  f[X\star_{\mathcal P}Y]
    \subseteq f[X]\diamond_{\mathcal P}f[Y].
\]

\medskip\noindent\textbf{Vom Produkt der Bilder zum Bild des Produkts.}
Sei umgekehrt \(z\in f[X]\diamond_{\mathcal P}f[Y]\). Dann gibt es
\(u\in f[X]\) und \(v\in f[Y]\) mit \(z=u\diamond v\).
Da \(u\) und \(v\) zu den jeweiligen Bildmengen gehören, können wir
\(x\in X\) und \(y\in Y\) wählen, sodass \(u=f(x)\) und \(v=f(y)\).
Nun gilt
\[
  z=u\diamond v=f(x)\diamond f(y)=f(x\star y).
\]
Das Element \(x\star y\) gehört zu \(X\star_{\mathcal P}Y\).
Also gehört \(z\) zu dessen Bild unter \(f\), und wir erhalten
\[
  f[X]\diamond_{\mathcal P}f[Y]
    \subseteq f[X\star_{\mathcal P}Y].
\]
Aus beiden Inklusionen folgt die behauptete Gleichheit. Dies ist der
Inhalt des ausgelagerten Resultats
\PowerIsoResultLink{SemigroupHomPreservesPowerProduct}{Ein Homomorphismus erhält das Mengenprodukt}.
Die formale Ableitung führt dieselben Schritte mit den einzelnen
Produkt- und Bildzeugen aus
(\PowerIsoProofLink{SemigroupHomPreservesPowerProduct}).

Für diesen Teil des Beweises brauchen wir weder Injektivität noch
Surjektivität. Die gewählten Urbilder müssen nicht eindeutig sein.
Ihre Existenz folgt schon daraus, dass wir mit Elementen der
Bildmengen \(f[X]\) und \(f[Y]\) begonnen haben. Insbesondere wird
keine Aussage darüber benötigt, ob jedes Element von \(B\) ein
Urbild unter \(f\) besitzt.

\clearpage
\section*{Die induzierte Potenzabbildung}
\hypertarget{power-isomorphisms.reading.operation}{}

Aus einer Funktion \(f\colon A\to B\) entsteht die Abbildung
\[
  \PowMapNE{f}\colon\PowNE{A}\longrightarrow\PowNE{B},
  \qquad X\longmapsto f[X].
\]
Wir nennen sie die von \(f\) induzierte Potenzabbildung.
Ihre Vorschrift lässt keine Wahl offen: Für jede Eingabemenge \(X\)
sind genau die Werte von \(f\) an den Elementen von \(X\) aufzunehmen.

Dass der angegebene Wertebereich stimmt, muss ebenfalls geprüft werden.
Aus \(X\subseteq A\) und \(f\colon A\to B\) folgt \(f[X]\subseteq B\).
Weil \(X\ne\varnothing\), gibt es ein \(x\in X\); dann gehört
\(f(x)\) zu \(f[X]\). Somit ist auch \(f[X]\ne\varnothing\).
Die Nichtleerheit geht also nicht verloren. Das entspricht
\FormulaRefAuto{NonemptyDirectImage}[thm].

Sei \(f\) wieder ein Homomorphismus. Auch
\(X\star_{\mathcal P}Y\) gehört zu \(\PowNE{A}\), sodass wir die
Potenzabbildung auf dieses Produkt anwenden dürfen. Mit der soeben
bewiesenen Gleichheit erhalten wir
\[
  \begin{aligned}
    \PowMapNE{f}(X\star_{\mathcal P}Y)
      &=f[X\star_{\mathcal P}Y]\\
      &=f[X]\diamond_{\mathcal P}f[Y]\\
      &=\PowMapNE{f}(X)\diamond_{\mathcal P}\PowMapNE{f}(Y).
  \end{aligned}
\]
Die erste und die letzte Gleichheit setzen nur die Vorschrift
\(\PowMapNE{f}(X)=f[X]\) ein
(\FormulaRefAuto{NonemptyPowerMapValue}[thm]). Der mittlere Schritt
enthält die eigentliche mathematische Aussage.

Dies ist das zweite ausgelagerte Resultat,
\PowerIsoResultLink{PowerSemigroupInducedMapOperation}{Operationserhaltung der induzierten Potenzabbildung}.
Die formale Ableitung protokolliert die Einsetzungen und die dafür
nötige Abgeschlossenheit
(\PowerIsoProofLink{PowerSemigroupInducedMapOperation}).
Da Ausgangs- und Zielstruktur bereits Halbgruppen sind, ist
\(\PowMapNE{f}\) damit ein Halbgruppenhomomorphismus.

\medskip\noindent\textbf{Was auf den Einermengen geschieht.}
Für jedes \(x\in A\) gilt
\[
  \PowMapNE{f}(\{x\})=\{f(x)\}.
\]
Auf den Einermengen erkennen wir daher die ursprüngliche Abbildung
wieder. Zugleich stimmt dort das Mengenprodukt mit dem ursprünglichen
Produkt überein:
\[
  \{x\}\star_{\mathcal P}\{y\}=\{x\star y\}.
\]
Die eben bewiesene Operationserhaltung setzt also das bekannte
Homomorphiegesetz von einzelnen Elementen auf beliebige nichtleere
Teilmengen fort. Wenn \(f\) verschiedene Elemente zusammenfasst,
können dabei auch verschiedene Teilmengen dasselbe Bild besitzen.
Für einen Isomorphismus werden wir nun zeigen, dass dies ausgeschlossen ist.

\clearpage
\section*{Warum aus einem Isomorphismus wieder ein Isomorphismus wird}
\hypertarget{power-isomorphisms.reading.isomorphism}{}

Sei jetzt \(f\colon A\to B\) ein Halbgruppenisomorphismus. Die
Operationserhaltung von \(\PowMapNE{f}\) ist bereits bewiesen.
Es bleibt die Bijektivität. Dafür geben wir eine Umkehrabbildung an.

Da \(f\) bijektiv ist, existiert die Umkehrfunktion
\(g=f^{-1}\colon B\to A\). Sie erfüllt
\[
  g(f(x))=x\quad(x\in A),
  \qquad f(g(y))=y\quad(y\in B).
\]
Für \(Y\in\PowNE{B}\) setzen wir
\[
  G(Y)=g[Y]=\{g(y)\mid y\in Y\}.
\]
Diese Menge ist eine nichtleere Teilmenge von \(A\). Die Vorschrift
definiert deshalb eine Funktion
\(G\colon\PowNE{B}\to\PowNE{A}\).

Wenn wir dafür \(f^{-1}[Y]\) schreiben, ist hier ausdrücklich das
\emph{Bild von \(Y\) unter der Umkehrfunktion} \(f^{-1}\) gemeint.
Bei einer bijektiven Funktion ist diese Menge zugleich das Urbild
\(\{x\in A\mid f(x)\in Y\}\). Die Umkehrfunktion steht uns hier
zur Verfügung, weil \(f\) bijektiv ist.

Für \(X\in\PowNE{A}\) gilt nun
\[
  G(\PowMapNE{f}(X))
    =g[f[X]]
    =\{g(f(x))\mid x\in X\}
    =X.
\]
Jedes Element von \(X\) wird erst mit \(f\) und anschließend mit
\(g\) abgebildet und kommt dabei zu sich selbst zurück. In der
anderen Richtung erhalten wir für \(Y\in\PowNE{B}\)
\[
  \PowMapNE{f}(G(Y))
    =f[g[Y]]
    =\{f(g(y))\mid y\in Y\}
    =Y.
\]
Damit ist \(G\) eine beidseitige Umkehrabbildung von
\(\PowMapNE{f}\). Insbesondere ist \(\PowMapNE{f}\) injektiv:
Aus gleichen Bildmengen erhalten wir durch Anwenden von \(G\)
die Gleichheit der Ausgangsmengen. Sie ist auch surjektiv:
Zu jeder nichtleeren Teilmenge \(Y\subseteq B\) ist \(G(Y)\)
eine Eingabemenge mit Bild \(Y\). Dies veranschaulicht den in der
formalen Ableitung verwendeten Satz
\FormulaRefAuto{NonemptyPowerMapBijective}[thm].

\medskip\noindent\textbf{Hauptsatz.}
Ist \(f\) ein Isomorphismus von \((A,\star)\) auf \((B,\diamond)\),
so ist
\[
  \PowMapNE{f}\colon
  (\PowNE{A},\star_{\mathcal P})
    \longrightarrow(\PowNE{B},\diamond_{\mathcal P}),
  \qquad X\longmapsto f[X],
\]
ein Isomorphismus der Potenzhalbgruppen. Seine Umkehrabbildung ist
\(Y\mapsto f^{-1}[Y]\). Dies ist
\FormulaRefAuto{SemigroupIsoInducesPowerSemigroupIso}[thm] in Band~28.
Der Schlussbeweis verbindet die Bijektivität mit der bereits
nachgewiesenen Operationserhaltung
(\PowerIsoProofLink{SemigroupIsoInducesPowerSemigroupIso}).

\clearpage
\section*{Ein vollständiges Beispiel mit drei Elementen}

Wir betrachten \(A=\{0,1,2\}\) mit der Addition modulo drei als
Operation \(\star\). Für die Umbenennung nehmen wir
\(B=\{a,b,c\}\) und setzen
\[
  f(0)=a,\qquad f(1)=b,\qquad f(2)=c.
\]
Die Operation \(\diamond\) auf \(B\) wird durch die rechte der
folgenden Tabellen festgelegt. Links steht die Tabelle von \(A\):
\[
  \begin{array}{c|ccc}
    \star&0&1&2\\\hline
    0&0&1&2\\
    1&1&2&0\\
    2&2&0&1
  \end{array}
  \qquad\qquad
  \begin{array}{c|ccc}
    \diamond&a&b&c\\\hline
    a&a&b&c\\
    b&b&c&a\\
    c&c&a&b
  \end{array}
\]
Der Zeilenfaktor steht jeweils links, der Spaltenfaktor rechts.
Jeder Eintrag der rechten Tabelle entsteht aus dem entsprechenden
Eintrag der linken durch \(f\). Somit gilt für alle neun Paare
\(f(x\star y)=f(x)\diamond f(y)\).

Die Addition modulo drei ist assoziativ. Auch \(\diamond\) ist
assoziativ: Schreiben wir drei Elemente von \(B\) als \(f(x),f(y),f(z)\),
so ergeben die beiden Klammerungen die Werte
\(f((x\star y)\star z)\) und \(f(x\star(y\star z))\), die gleich sind.
Die Bijektivität von \(f\) ist an den drei verschiedenen Werten
\(a,b,c\) erkennbar. Damit liegt tatsächlich ein Isomorphismus von
Halbgruppen vor.

Die beiden Potenzhalbgruppen haben je sieben Elemente. Die induzierte
Abbildung ordnet sie vollständig wie folgt einander zu:
\[
  \begin{array}{c|c|c}
    \text{Kürzel}&X\subseteq A&\PowMapNE{f}(X)\subseteq B\\\hline
    E_0&\{0\}&\{a\}\\
    E_1&\{1\}&\{b\}\\
    E_2&\{2\}&\{c\}\\
    D_0&\{1,2\}&\{b,c\}\\
    D_1&\{0,2\}&\{a,c\}\\
    D_2&\{0,1\}&\{a,b\}\\
    T&A&B
  \end{array}
\]
Die Kürzel bezeichnen jeweils die entsprechende Menge auf der
betrachteten Seite. \(D_i\) ist auf der linken Seite die Menge,
in der genau \(i\) fehlt. Auf der rechten Seite fehlt entsprechend
\(f(i)\). Die sieben verschiedenen Bildmengen bilden die gesamte
nichtleere Potenzmenge von \(B\). Die Umkehrabbildung liest dieselbe
Tabelle von rechts nach links.

\clearpage
\section*{Die Mengenprodukte im Beispiel}

Mit den eben eingeführten Kürzeln lautet die vollständige
Multiplikationstabelle beider Potenzhalbgruppen:
\[
  \begin{array}{c|ccccccc}
    &E_0&E_1&E_2&D_0&D_1&D_2&T\\\hline
    E_0&E_0&E_1&E_2&D_0&D_1&D_2&T\\
    E_1&E_1&E_2&E_0&D_1&D_2&D_0&T\\
    E_2&E_2&E_0&E_1&D_2&D_0&D_1&T\\
    D_0&D_0&D_1&D_2&T&T&T&T\\
    D_1&D_1&D_2&D_0&T&T&T&T\\
    D_2&D_2&D_0&D_1&T&T&T&T\\
    T&T&T&T&T&T&T&T
  \end{array}
\]
Die Tabelle lässt sich in wenigen Schritten vollständig berechnen.
Für Einermengen gilt \(E_iE_j=E_{i+j}\), wobei die Indizes modulo drei
gelesen werden. Das Addieren eines festen Elements \(i\) verschiebt
die beiden Elemente von \(D_j\); genau \(i+j\) fehlt anschließend.
Also gilt \(E_iD_j=D_jE_i=D_{i+j}\).

Zwei zweielementige Teilmengen ergeben stets ganz \(A\). Um das für
jeden gewünschten Wert \(z\in A\) zu sehen, betrachten wir die beiden
zweielementigen Mengen \(D_i\) und
\(z-D_j=\{z-y\mid y\in D_j\}\), wobei auch die Subtraktion
modulo drei gerechnet wird. Sie müssen sich in der
Dreiermenge \(A\) schneiden. Für ein gemeinsames Element \(x\)
gibt es daher ein \(y\in D_j\) mit \(x=z-y\), also \(x\star y=z\).
Somit gilt \(D_iD_j=T\). Schließlich ergibt das Produkt von \(A\)
mit jeder nichtleeren Teilmenge wieder \(A\): Schon die Addition
eines festen Elements permutiert alle drei Elemente von \(A\).
Dies erklärt sämtliche Einträge der Tabelle. Durch die Umbenennung
gelten dieselben Rechnungen auf der rechten Seite.

Ein konkreter Eintrag zeigt das Zusammenspiel beider Wege. Für
\(X=\{0,1\}=D_2\) und \(Y=\{1,2\}=D_0\) ist
\[
  X\star_{\mathcal P}Y
    =\{0\star1,0\star2,1\star1,1\star2\}
    =\{1,2,2,0\}=A.
\]
Daher ist \(f[X\star_{\mathcal P}Y]=\{a,b,c\}\).
Berechnen wir zuerst die Bildmengen, so erhalten wir ebenfalls
\[
  \begin{aligned}
    f[X]\diamond_{\mathcal P}f[Y]
      &=\{a,b\}\diamond_{\mathcal P}\{b,c\}\\
      &=\{a\diamond b,a\diamond c,b\diamond b,b\diamond c\}\\
      &=\{b,c,c,a\}=\{a,b,c\}.
  \end{aligned}
\]
Die zweimal auftretenden Werte \(2\) beziehungsweise \(c\) werden
jeweils zu einem Mengenelement zusammengefasst.

\clearpage
\section*{Die anschließende Umkehrfrage}

Der Satz liefert eine klare Richtung: Ein Isomorphismus der
Grundhalbgruppen erzeugt einen Isomorphismus ihrer Potenzhalbgruppen.
Er beantwortet noch nicht die Rekonstruktionsfrage, ob aus der
Isomorphie zweier Potenzhalbgruppen auf die Isomorphie der
Grundhalbgruppen geschlossen werden kann.

Auch die Frage, ob ein vorgegebener Isomorphismus von Potenzhalbgruppen
durch eine Abbildung der Grundmengen entsteht, ist damit nicht entschieden.
Die hier konstruierten Abbildungen haben die besondere Form
\(X\mapsto f[X]\) und senden insbesondere Einermengen auf Einermengen.
Für einen beliebig vorgegebenen Isomorphismus der Potenzhalbgruppen
müsste man solche Eigenschaften erst begründen. Die bewiesene
Übertragung stellt dafür den Ausgangspunkt bereit.

\medskip\noindent\PowerIsoMainLink{}\enspace\textperiodcentered\enspace\PowerIsoProofsLink{}


